%% P11 --- SVD, low-rank approximation and conditioning
%%
%% Every number in this program that the reader cannot do in their head comes
%% from code/p11_svd_conditioning.py and is pulled in with \val{}. The console
%% listing is a committed transcript written by that script.
%%
%% THE THESIS: every matrix, square or not, factors into rotate-stretch-rotate;
%% truncating the stretch gives the best approximation of a given rank; and the
%% ratio of largest to smallest stretch is the number that says how much error
%% a solve will amplify.
%%
%% WHAT THIS PROGRAM IS OWED, read out of the files rather than remembered:
%%   P10  HANDS IT OVER BY NAME -- "every matrix, of any shape, turns out to be
%%        a rotation, then a stretch along perpendicular axes, then another
%%        rotation, so the picture survives even where the theorem does not."
%%        P10's THREE FAILURES are this program's motivation, and its QDQ^T
%%        construction is reused with U and V allowed to differ.
%%   P09  hands over the accuracy half of "do not invert" twice and by name:
%%        the operation count was measured there, "the accuracy is P11's, and
%%        it is the half that decides the argument". Paid in section 5.
%%   P08  names a debt in a warning box: rank collapse "needs Program P11's
%%        machinery and a real model". The machinery is here. The model is not,
%%        and section 3 says so rather than manufacturing one.
%%   P05  gives norms and projection; P04 the basis; P02 what a lost digit is.
%%
%% WHAT THIS PROGRAM LEAVES ALONE, checked against tools/programs.json:
%%   the Hessian and the step-size bound                              -> P17
%%   the optimisers and how many iterations conditioning costs        -> P20
%%
%% FORWARD REFERENCE: the covariance matrix is P24's. Program P10 DISCHARGED
%% it by proving both facts it needs locally, so this program cites P10 rather
%% than re-declaring a debt that has been paid.
%%
%% Twin: programs/pl/P11-svd-conditioning.tex. Frame for frame, macro for
%% macro, number for number; tools/parity.py compares them.

\program{SVD, low-rank approximation and conditioning}
\label{prog:P11}
\index{singular value decomposition}
\index{matrix!condition number}

\begin{outcomes}
\outcome{1--6}{say what the three factors of a singular value decomposition do,
  and why every matrix has one}
\outcome{7--11}{say why a singular value is not an eigenvalue, and what a
  rectangular matrix has instead}
\outcome{12--19}{say what a rank-$r$ truncation keeps and what it costs, and
  what a spectrum has to look like for it to be cheap}
\outcome{20--24}{read the condition number as an amplification factor from
  input error to output error}
\outcome{25--33}{say why forming $A\T A$ is the wrong way to solve a
  least-squares problem, with the digits}
\end{outcomes}

\begin{quiz}
\item A matrix takes a vector from a space of one size to a space of another.
  Can it have an eigenvector? \teachesat{1--2}
  \answerto{No \dash{} $Av = \lambda v$ needs $Av$ and $v$ to live in the same
  space, so the equation cannot even be written. Singular values are defined
  for every matrix of every shape.}
\item $A$ has singular values $\val{p11.sig.hi}$ and $\val{p11.sig.lo}$. Give
  the eigenvalues of $A\T A$. \teachesat{7--8}
  \answerto{$100$ and $4$ \dash{} the squares. $A\T A = V\Sigma^{2}V\T$, so its
  eigenvalues are the squared singular values whatever $A$ was.}
\item A rank-$1$ approximation is subtracted from $A$. What is left, in terms
  of the singular values? \teachesat{14--15}
  \answerto{Exactly $\sigma_2$, if the piece removed was the largest one. Each
  rank-one piece contributes its own singular value and nothing else does.}
\item $A$ has singular values between $\val{p11.sig.lo}$ and
  $\val{p11.sig.hi}$. By how much can it amplify a relative error, worst case?
  \teachesat{20--21}
  \answerto{$\val{p11.kappa}$ times \dash{} the ratio of the largest stretch to
  the smallest. That ratio is the condition number.}
\item You form $A\T A$ to solve a least-squares problem. What have you done to
  the condition number? \teachesat{25--26}
  \answerto{Squared it. $A\T A$ has the squared singular values, so its
  condition number is $\val{p11.kappa.sq}$ where $A$'s was
  $\val{p11.kappa}$ \dash{} and that costs digits.}
\end{quiz}

% ============================================================
\section{What survives when the theorem does not}
\label{sec:P11-svd}
\index{singular value decomposition!definition}

\begin{fr}
Program~\ref{prog:P10} ended with three matrices it could not handle, and they
were not exotic ones.

A non-symmetric matrix whose eigenvectors met at
$\val{p10.nonsym.angle}$ degrees rather than a right angle. A quarter turn with
no real eigenvalue at all. A shear with one eigenvalue twice and only one
direction to show for it.

And a fourth case it never even raised, which is the one an engineer meets
most. Write down the eigenvector equation $Av = \lambda v$ for a matrix with
more columns than rows, and say what goes wrong.
\dotline
\nextframe
\end{fr}

\begin{fr}
\begin{ansblock}
$Av$ and $v$ live in spaces of different sizes, so the equation cannot be
written down at all.
\end{ansblock}

An eigenvector is a direction that comes back along itself, and \enquote{along
itself} requires the input and the output to be in the same room. A
$\val{p11.rect.rows} \times \val{p11.rect.cols}$ matrix takes vectors from a
space of three dimensions to one of two. Nothing it produces can be compared
with what it consumed.

So the whole of Program~\ref{prog:P10} applies to square matrices, and most of
the matrices in a network are not square.
\end{fr}

\begin{fr}
The repair keeps the picture and gives up the constraint.

Program~\ref{prog:P10} asked for \emph{one} set of perpendicular directions
that the matrix merely scales. That is asking a lot: the same directions have
to work at both ends.

Ask instead for \textbf{two} sets \dash{} perpendicular directions in the input
space, and perpendicular directions in the output space, with the matrix
carrying the $i$th of the first to a multiple of the $i$th of the second:
\[ Av_i = \sigma_i u_i \]

Every matrix has that, of every shape, with no exceptions and no conditions.
The $\sigma_i$ are its \textbf{singular values}.
\end{fr}

\begin{fr}
Written as a factorisation, with the $v$'s as the columns of $V$ and the $u$'s
as the columns of $U$:
\[ A = U\Sigma V\T \]
where $U$ and $V$ are orthogonal \dash{} Program~\ref{prog:P09}'s word: their
transposes are their inverses \dash{} and $\Sigma$ is diagonal with the
singular values on it, conventionally largest first.

Read right to left, that is three moves. $V\T$ turns the input space so the
special directions lie along the axes. $\Sigma$ stretches each axis by its own
number. $U$ turns the result into wherever the output lives.

\textbf{Rotate, stretch, rotate.} Every matrix, without exception.
\end{fr}

\mermaidfig{p11-three-moves}{The only three things any matrix does, in the
order it does them.}{the SVD as three moves}

\begin{fr}
Program~\ref{prog:P10} checked its theorem exactly by building the answer
rather than solving for it, and the same trick works here with one change worth
noticing.

There, $A = QDQ\T$ used the \emph{same} $Q$ on both sides, because a symmetric
matrix's two sets of directions are the same set. Here they need not be, so the
construction uses two different rational rotations \dash{} one from the
$3$--$4$--$5$ triangle and one from $5$--$12$--$13$ \dash{} and forms
$A = U\Sigma V\T$ with $\Sigma$ carrying $\val{p11.sig.hi}$ and
$\val{p11.sig.lo}$.

The resulting $A$ has thoroughly unpleasant entries, $\frac{246}{65}$ among
them, and is not symmetric. Its singular values are $\val{p11.sig.hi}$ and
$\val{p11.sig.lo}$ exactly, because they were put there, and the script checks
$Av_i = \sigma_i u_i$ for both $i$ over fractions.

Before reading on, predict one thing about it. Given only that its singular
values are $\val{p11.sig.hi}$ and $\val{p11.sig.lo}$, what is
$\lvert\det A\rvert$?
\dotline
\nextframe
\end{fr}

\begin{fr}
\begin{ansblock}
$\val{p11.detA}$, which is
$\val{p11.sig.hi} \times \val{p11.sig.lo}$.
\end{ansblock}

That is not a coincidence and it is not new. $\det(U\Sigma V\T) = \det U \cdot
\det\Sigma \cdot \det V\T$ by Program~\ref{prog:P09}'s product rule, the two
rotations have determinant $\pm 1$, and $\det\Sigma$ is the product of the
singular values. \textbf{The determinant is the product of the stretches, with
the rotations contributing nothing} \dash{} which is what Program~\ref{prog:P09}
meant by volume, said in this program's vocabulary.
\end{fr}

% ============================================================
\section{Not eigenvalues}
\label{sec:P11-not-eigen}
\index{singular value decomposition!against eigenvalues}

\begin{fr}
The two ideas sit close enough together to be confused constantly, so here is
the relationship, exactly.

Take $A = U\Sigma V\T$ and form $A\T A$. Multiply it out, using
$U\T U = I$.
\dotline
\nextframe
\end{fr}

\begin{fr}
\begin{ansblock}
$A\T A = V\Sigma\T U\T U \Sigma V\T = V\Sigma^{2}V\T$.
\end{ansblock}

Which is Program~\ref{prog:P10}'s $QDQ\T$ exactly, with $V$ for $Q$ and
$\Sigma^{2}$ for $D$.

So: \textbf{the singular values of $A$ are the square roots of the eigenvalues
of $A\T A$}, and the right singular vectors are that matrix's eigenvectors.
$A\T A$ is symmetric whatever $A$ was, so the spectral theorem always applies
to it, and every $\sigma_i$ is real and non-negative because it is the square
root of a non-negative number.

That is why singular values exist for every matrix: they are borrowed from a
symmetric matrix you can always build.
\end{fr}

\begin{fr}
Now the trap, and this example is chosen to make it hard rather than easy.

The constructed $A$ has singular values $\val{p11.sig.hi}$ and
$\val{p11.sig.lo}$. It is square, so it also has eigenvalues of its own.

Write down what you would guess they are, then read on.
\dotline
\nextframe
\end{fr}

\begin{fr}
\begin{ansblock}
$\val{p11.lam.hi}$ and $\val{p11.lam.lo}$. Not $\val{p11.sig.hi}$ and
$\val{p11.sig.lo}$.
\end{ansblock}

Close. Close enough that on a plot you would not see the difference, and close
enough that a table quoting the wrong pair would pass a review.

\transcript{p11-singular-vs-eigen}

The listing takes the other route: the trace and determinant of $A\T A$ are
$104$ and $400$, which are the sum and product of $100$ and $4$ \dash{} the
squared singular values, exactly.
\end{fr}

\begin{trapbox}
\textbf{\enquote{Singular values are eigenvalues.}}

They coincide for one family: symmetric positive semi-definite matrices, where
$A\T A = A^{2}$ and the square root undoes the square. That family includes
covariance matrices and Gram matrices, which is why the habit forms.

Everywhere else they differ, and the two differ in kind and not only in value:

\begin{itemize}
\item singular values exist for \textbf{every} matrix, including rectangular
  ones; eigenvalues need a square one;
\item singular values are \textbf{real and non-negative}, always; eigenvalues
  can be complex, as Program~\ref{prog:P10}'s quarter turn showed;
\item there are always \textbf{as many singular values as the smaller
  dimension}, with a full set of directions; eigenvectors can run short, as the
  shear showed.
\end{itemize}

The example above is the uncomfortable case: $\val{p11.lam.hi}$ against
$\val{p11.sig.hi}$ is a $5\%$ difference, and a plot of one labelled as the
other would look entirely reasonable.
\end{trapbox}

\begin{fr}
For a rectangular matrix there are two Gram matrices rather than one, and the
sizes are worth seeing.

A $\val{p11.rect.rows} \times \val{p11.rect.cols}$ matrix $R$ gives
$R\T R$ of size $\val{p11.rect.cols} \times \val{p11.rect.cols}$ and $RR\T$ of
size $\val{p11.rect.rows} \times \val{p11.rect.rows}$. Different matrices,
different sizes, and neither of them is $R$.

Their traces agree at $\val{p11.rect.energy}$, which the script checks, and
that is the tell: the two carry the same total, so the singular values they
imply are the same list. \textbf{The larger Gram simply has extra zeros}, one
for each dimension the matrix cannot reach.
\end{fr}

% ============================================================
\section{Keeping the largest pieces}
\label{sec:P11-lowrank}
\index{matrix!low-rank approximation}

\begin{fr}
The factorisation can be read a second way, and the second reading is the
useful one.

$U\Sigma V\T$ multiplies out to a \emph{sum}:
\[ A = \sigma_1 u_1v_1\T + \sigma_2u_2v_2\T + \cdots \]
Each term is a column times a row, which by Program~\ref{prog:P08} is a matrix
of rank exactly one. So a matrix is a sum of rank-one pieces, one for each
singular value, weighted by it.

The pieces do not interact: the $u$'s are perpendicular and so are the $v$'s.
Say what that buys \dash{} what does the total \enquote{size} of $A$ decompose
into?
\dotline
\nextframe
\end{fr}

\begin{fr}
\begin{ansblock}
The squares of the singular values, and nothing else.
\end{ansblock}

The sum of the squares of all the entries of $A$ \dash{} the quantity
Program~\ref{prog:P05} would call the squared length of $A$ read as one long
vector \dash{} is $\sigma_1^{2} + \sigma_2^{2} + \cdots$.

Perpendicular pieces add in squares, which is Program~\ref{prog:F09}'s
Pythagoras with matrices in place of arrows. So each singular value carries a
share of the whole, and the share is $\sigma_i^{2}$.
\end{fr}

\mermaidfig{p11-keep-the-largest}{Why the pieces of a matrix can be weighed
one at a time.}{a matrix as a sum of pieces}

\begin{fr}
Now the approximation. Keep the first $r$ terms and throw the rest away.

You have just been told what each piece contributes. So without computing
anything: what is left over, measured the same way?
\dotline
\nextframe
\end{fr}

\begin{fr}
\begin{ansblock}
$\sigma_{r+1}^{2} + \sigma_{r+2}^{2} + \cdots$ \dash{} exactly the pieces you
discarded.
\end{ansblock}

On the constructed matrix of section 1, truncating to rank $1$ leaves an error
of exactly $\val{p11.ey.err}$, which is $\sigma_2$, and the script checks that
over fractions.

The error is not \emph{about} $\sigma_2$ or bounded by it. It is $\sigma_2$.

Which raises the question the rest of the section is about. Truncating is
\emph{one} way to build a rank-$1$ matrix near $A$. Is there a better one?
\dotline
\nextframe
\end{fr}

\begin{fr}
\begin{ansblock}
No. There is no rank-$1$ matrix closer to $A$ than that one.
\end{ansblock}

\textbf{Eckart--Young}, which this book states and does not prove. Of all
matrices of rank $r$, the truncated SVD is the closest to $A$. Not a good
choice, not a defensible heuristic: the best one there is.

That is a strong claim \dash{} it quantifies over every rank-$r$ matrix, of
which there are infinitely many \dash{} so it deserves the same treatment
Program~\ref{prog:P10} gave the spectral theorem. The script builds
$\val{p11.ey.trials}$ other rank-one matrices from rational direction pairs,
takes the best scaling for each, and requires every one to be at least as bad.
Every one is.
\end{fr}

\begin{fr}
As in Program~\ref{prog:P10}, those trials are evidence about the code rather
than about the theorem: a counterexample would have refuted a proof.

What they are for is different here, and worth separating. The spectral theorem
was checked because the word \emph{exactly} needed exact arithmetic.
Eckart--Young is checked because \textbf{the claim is easy to misremember as
weaker than it is} \dash{} \emph{a good approximation} rather than \emph{the
best one} \dash{} and seeing thousands of alternatives all lose is what fixes
which of the two the reader carries away.
\end{fr}

\begin{fr}
Here is what that buys, and it is the reason anybody outside a mathematics
department cares.

Take a matrix whose singular values run $100, 40, 16, 6, 2, 1$. Keep the top
$\val{p11.spec.keep}$ of the $\val{p11.spec.n}$ and throw four away.

What fraction of the total do you still have?
\dotline
\nextframe
\end{fr}

\begin{fr}
\begin{ansblock}
$\val{p11.spec.pct}\%$, leaving $\val{p11.spec.rest.pct}\%$ behind.
\end{ansblock}

A third of the pieces carrying almost all of the content. That is the whole
argument for storing $U_r\Sigma_r V_r\T$ instead of $A$, and for
Program~\ref{prog:P08}'s low-rank adapter: if a matrix is nearly rank
$\val{p11.spec.keep}$, treating it as rank $\val{p11.spec.keep}$ costs almost
nothing.

Notice the shape of the condition. It is not about the size of the matrix. It
is about how fast the singular values \textbf{decay}, and a matrix whose
spectrum is flat gets no benefit at all.
\end{fr}

\begin{warning}
\textbf{That spectrum was constructed, and this book has not measured a real
one.}

The frame above chose six numbers that decay quickly and then reported what
follows. It demonstrates the \emph{mechanism} and it is not evidence that any
particular matrix behaves that way.

The empirical claim \dash{} that a trained model's embedding or weight matrices
have spectra concentrated enough for a rank-$r$ approximation to be nearly
free \dash{} is what the low-rank adaptation literature rests on, and settling
it needs the singular values of a real model's real matrices. This book does
not have one and does not download one.

It is the same debt Program~\ref{prog:P08} declined to pretend it had paid for
rank collapse, and the same answer: the machinery is here, the model is not,
and saying which is which is the whole of the discipline.
\end{warning}

% ============================================================
\section{Error in, error out}
\label{sec:P11-conditioning}
\index{matrix!condition number}

\begin{fr}
A different question about the same three moves, and it is the one that decides
whether an answer can be trusted.

An input arrives already carrying error \dash{} from a measurement, from a
sensor, or simply from the last rounding Program~\ref{prog:P01} described.
The matrix acts on the error exactly as it acts on everything else.

If the input error happens to point along $v_1$ it is multiplied by
$\sigma_1$. If it points along the last $v$ it is multiplied by the smallest
$\sigma$. The question is what happens to the error \emph{relative} to the
answer, and that is a ratio of two of those.
\dotline
\nextframe
\end{fr}

\begin{fr}
\begin{ansblock}
Worst case, the relative error is multiplied by
$\sigma_{\max}/\sigma_{\min}$.
\end{ansblock}

The worst case has the answer as small as the map allows and the error as large
as it allows, so the two extremes divide.

That ratio is the \textbf{condition number}, written $\kappa(A)$, and for
section 1's matrix it is
$\val{p11.sig.hi}/\val{p11.sig.lo} = \val{p11.kappa}$.

It is dimensionless, and it multiplies \emph{relative} error. So: an input good
to sixteen significant figures goes into a matrix with $\kappa = 10^{6}$.
Roughly how many figures is the answer good to?
\dotline
\nextframe
\end{fr}

\mermaidfig{p11-error-in-error-out}{What a matrix does to error that arrived
with the input.}{the condition number as amplification}

\begin{fr}
\begin{ansblock}
About ten. Six of the sixteen have gone.
\end{ansblock}

The rule of thumb is the one to carry: \textbf{you lose about
$\logb{10}\kappa$ decimal digits}. Double precision starts with about sixteen,
so a condition number of $10^{16}$ leaves nothing at all, and one of
$\val{p11.spec.kappa}$ costs two digits of the sixteen and is unremarkable.
\end{fr}

\begin{fr}
Two things it is not, and both are common readings.

It is not a property of the \emph{algorithm}. A condition number is a fact
about the matrix, and it says what any correct method must contend with. A
badly conditioned problem solved by a perfect algorithm still gives a poor
answer, because the answer really is that sensitive to the input.

And it is not a bug report. Program~\ref{prog:P10}'s ravine had a large
condition number by design; so does any problem where one direction matters far
more than another. The number tells you how much precision the problem will
cost, not that somebody has made a mistake.

Both readings come from the same place: a number that predicts a bad outcome
gets treated as a diagnosis. It is closer to a weather forecast. It tells you
what the conditions are and how much they will cost you, and the honest
response is to budget for the loss rather than to look for the mistake that
caused it.

Which is worth joining to something already measured. For a symmetric positive
definite matrix the singular values \emph{are} the eigenvalues. So what has
Program~\ref{prog:P10} already called $\kappa$, without using the word?
\dotline
\nextframe
\end{fr}

\begin{fr}
\begin{ansblock}
The ratio of the eigenvalues of a basin \dash{} whose square root was how
elongated the level ellipse is.
\end{ansblock}

That gives the connection two programs have been circling. A ravine is a basin
with a large $\kappa$; a large $\kappa$ is a problem that gives up digits; and
Program~\ref{prog:P20} will show it is also a problem that costs an optimiser
iterations.

\textbf{One number, three consequences}, and this program's contribution is the
middle one.
\end{fr}

% ============================================================
\section{Why the code does not form \texorpdfstring{$A\T A$}{A transpose A}}
\label{sec:P11-normal}
\index{least squares!normal equations}

\begin{fr}
Program~\ref{prog:P09} measured that solving a system costs about a fifth of
what inverting it costs, and then said something it could not back up:

\begin{quote}
The operation count is measured here. The accuracy is Program~\ref{prog:P11}'s,
and it is the half that decides the argument.
\end{quote}

This section pays that. The vehicle is least squares, because it is where the
temptation is strongest: Program~\ref{prog:P08} derived the normal equations
$A\T A\hat{x} = A\T b$ as the closed form, and a closed form asks to be typed
in as written.

You know what the singular values of $A\T A$ are. So what is its condition
number?
\dotline
\nextframe
\end{fr}

\begin{fr}
\begin{ansblock}
$\kappa(A)^{2}$. The singular values are squared, so their ratio is squared.
\end{ansblock}

For section 1's matrix that is $\val{p11.kappa.sq}$ against
$\val{p11.kappa}$, which the script checks exactly.

And now put it beside the rule of thumb two frames ago. If you lose about
$\logb{10}\kappa$ digits, and forming $A\T A$ squares $\kappa$, then forming
$A\T A$ \textbf{doubles the number of digits you lose}. Not a little worse:
twice the loss, for a step taken only because the formula was written that way.
\end{fr}

\begin{fr}
That is an argument. Here is the measurement, and it is the one claim in this
program that is genuinely about floating point rather than about matrices.

A real least-squares problem: fit a polynomial of degree $\val{p11.ls.deg}$
through $\val{p11.ls.points}$ points spread over $\intcc{0}{1}$. The
coefficients are chosen, so the exact answer is known, and the right-hand side
is built over fractions so that nothing is lost before the solving starts.

Then the same system is solved twice in ordinary double precision \dash{} once
through $A\T A$, and once by a factorisation that never forms it.
\end{fr}

\begin{fr}
\begin{tabularx}{\linewidth}{@{}lX@{}}
\toprule
\textbf{route} & \textbf{relative error in the coefficients} \\
\midrule
factorisation, $A\T A$ never formed & better than $\val{p11.ls.qr.bound}$ \\
the normal equations & worse than $\val{p11.ls.ne.floor}$ \\
\bottomrule
\end{tabularx}

\medskip
At least $\val{p11.ls.digits}$ decimal digits, given away for typing the
formula as written. Both figures are quoted as bounds rather than as
measurements, and the reason is Program~\ref{prog:P01}'s: a residual is a
property of the machine that produced it, so what is committed here is the gap,
which is not.
\end{fr}

\begin{fr}
The failure has the shape that makes it dangerous rather than the shape that
makes it obvious.

Nothing raises. No warning appears. The coefficients that come back look like
coefficients: right order of magnitude, right signs, plausible to anyone who
did not know the answer in advance. They are simply wrong from the sixth
significant figure onwards.

Program~\ref{prog:P07} named this shape exactly \dash{} a defect that is mild
while you are debugging and severe once things are working is the worst shape a
defect can have. Here the analogue is: a defect that looks like an answer is
worse than one that looks like a crash.
\end{fr}

\begin{fr}
So the advice \enquote{do not form the normal equations} is now two claims with
two owners, and both have numbers.

Program~\ref{prog:P09}: forming an inverse costs about
$\val{p09.cost.ratio}$ times the arithmetic of solving, measured by counting.
This program: forming $A\T A$ squares the condition number, so it doubles the
digits lost, measured by solving the same system twice.

The second is the half that decides it. A factor of $\val{p09.cost.ratio}$ in
speed is an engineering trade-off that a faster machine can absorb. Digits are
not recoverable at any speed.
\end{fr}

\begin{fr}
One more object, and it is a definition rather than a technique.

Program~\ref{prog:P08} showed a tall system usually has no solution and that
the best answer is a projection. The SVD gives that answer a closed form: turn
by $U\T$, divide by each singular value, turn by $V$. Written as a matrix it is
called the \textbf{pseudoinverse}, $A^{+}$.

Where $A$ is invertible it agrees with $A^{-1}$. Where it is not, it gives the
least-squares answer, and where several answers tie it gives the shortest.

Note that all three steps are things this program has already given names to.
Turning by an orthogonal matrix is Program~\ref{prog:P09}'s lossless change of
basis; dividing by the singular values is the only place any scaling happens;
and turning back is the same move again. Nothing in the definition is new,
which is why it is a definition rather than a technique.

That middle step is worth staring at. It divides by \emph{each} singular value,
including the smallest. Say what that does in a direction the matrix nearly
flattened.
\dotline
\nextframe
\end{fr}

\begin{fr}
\begin{ansblock}
It nearly explodes: a division by something nearly zero.
\end{ansblock}

That is the condition number again, arriving as a division you can point at.

It is also why truncating helps: dropping the smallest singular values before
inverting bounds the amplification, at the cost of not answering in those
directions at all. Whether that trade is worth taking is a judgement about the
problem, and this book does not have a rule for it.
\end{fr}

\begin{fr}
Part III closes here, and it closes on one factorisation.

Program~\ref{prog:P04} asked what a space is and Program~\ref{prog:P06} made a
matrix a function. Program~\ref{prog:P08} counted how much of the space
survives it, Program~\ref{prog:P09} measured how much, and
Program~\ref{prog:P10} found the directions that survive unturned \dash{} for
the matrices that have them.

$A = U\Sigma V\T$ has all of it. The rank is how many singular values are
non-zero; the determinant is their product; the null space is spanned by the
$v$'s whose $\sigma$ is zero; the best low-rank approximation is the first few
terms; and the condition number is the first divided by the last.
\textbf{Six programs' worth of questions, answered by reading one
factorisation}, which is the strongest reason to know it exists.
\end{fr}

% ============================================================
\begin{summarybox}
\sumitem{1--2}{\textbf{An eigenvector equation cannot be written for a
  rectangular matrix at all}, because $Av$ and $v$ live in spaces of different
  sizes. Program~\ref{prog:P10} therefore applies to square matrices, and most
  of the matrices in a network are not square.}
\sumitem{3--4}{The repair asks for \textbf{two} sets of perpendicular
  directions rather than one \dash{} $Av_i = \sigma_i u_i$ \dash{} and every
  matrix of every shape has them. Written as $A = U\Sigma V\T$ that is
  \textbf{rotate, stretch, rotate}, read right to left, with no exceptions and
  no conditions.}
\sumitem{5--6}{Checked exactly on a matrix built from \emph{two} different
  rational rotations, which is Program~\ref{prog:P10}'s construction with the
  constraint that $U = V$ lifted. And $\lvert\det A\rvert$ is the product of
  the singular values, because the rotations contribute nothing \dash{} which
  is Program~\ref{prog:P09}'s volume in this program's vocabulary.}
\sumitem{7--8}{$A\T A = V\Sigma^{2}V\T$, so \textbf{the singular values are the
  square roots of the eigenvalues of $A\T A$}. That is why they exist for every
  matrix: they are borrowed from a symmetric matrix you can always build, and
  the spectral theorem always applies to it.}
\sumitem{9--10}{\textbf{Singular values are not eigenvalues}, and the example
  is chosen to make the confusion understandable: $\val{p11.lam.hi}$ and
  $\val{p11.lam.lo}$ against $\val{p11.sig.hi}$ and $\val{p11.sig.lo}$, a
  $5\%$ difference that no plot would show. They differ in kind too \dash{}
  singular values exist for every shape, are always real and non-negative, and
  never run short of directions.}
\sumitem{11}{A rectangular matrix has \textbf{two} Gram matrices of different
  sizes, whose traces agree at $\val{p11.rect.energy}$. The larger simply
  carries extra zeros, one for each dimension the matrix cannot reach.}
\sumitem{12--13}{$A$ is a \textbf{sum of rank-one pieces}, one per singular
  value, and they do not interact \dash{} so the total size decomposes into
  $\sigma_i^{2}$, which is Program~\ref{prog:F09}'s Pythagoras with matrices in
  place of arrows.}
\sumitem{14--15}{Truncating to rank $r$ therefore leaves \textbf{exactly} the
  pieces discarded: $\val{p11.ey.err}$ for the worked matrix, which is
  $\sigma_2$. Not about $\sigma_2$, not bounded by it.}
\sumitem{16--17}{\textbf{Eckart--Young}: of all matrices of rank $r$ the
  truncation is the closest. Stated, not proved, and checked against
  $\val{p11.ey.trials}$ alternatives \dash{} not as evidence for the theorem
  but because \textbf{the claim is easy to misremember as weaker than it is},
  and watching thousands of alternatives lose fixes which version the reader
  carries away.}
\sumitem{18--19}{Two of six singular values carrying $\val{p11.spec.pct}\%$ is
  the entire argument for storing a truncation and for a low-rank adapter. And
  the condition is about \textbf{how fast the spectrum decays}, not about the
  size of the matrix: a flat spectrum gets no benefit at all.}
\sumitem{20--21}{Error arrives with the input and is stretched like everything
  else, so worst case the \emph{relative} error is multiplied by
  $\sigma_{\max}/\sigma_{\min}$. That is the \textbf{condition number}, and for
  the worked matrix it is $\val{p11.kappa}$.}
\sumitem{22--23}{\textbf{You lose about $\logb{10}\kappa$ decimal digits}, out
  of the sixteen a double starts with. It is not a property of the algorithm
  \dash{} a perfect method still gives a poor answer on a badly conditioned
  problem \dash{} and it is not a bug report.}
\sumitem{24}{For a symmetric positive definite matrix $\kappa$ is the
  eigenvalue ratio Program~\ref{prog:P10} measured, so \textbf{a ravine, a loss
  of digits and a slow optimiser are one number read three ways}. The middle
  one is this program's.}
\sumitem{25--26}{$A\T A$ has the \emph{squared} singular values, so
  \textbf{$\kappa(A\T A) = \kappa(A)^{2}$} \dash{} $\val{p11.kappa.sq}$ against
  $\val{p11.kappa}$. Squaring $\kappa$ doubles the digits lost, for a step
  taken only because Program~\ref{prog:P08}'s closed form was written that
  way.}
\sumitem{27--29}{Measured on a degree-$\val{p11.ls.deg}$ fit through
  $\val{p11.ls.points}$ points: the factorisation clears
  $\val{p11.ls.qr.bound}$, the normal equations fail $\val{p11.ls.ne.floor}$,
  \textbf{at least $\val{p11.ls.digits}$ digits given away}. Committed as
  bounds, because a residual belongs to the machine. And nothing raises: the
  wrong coefficients look like coefficients, which is worse than a crash.}
\sumitem{30}{So \enquote{do not form the normal equations} is two claims with
  two owners. Program~\ref{prog:P09} measured the arithmetic, this program the
  accuracy \dash{} and \textbf{the accuracy is the half that decides it},
  because a factor in speed is absorbed by a faster machine and digits are not
  recoverable at any speed.}
\sumitem{31--32}{The \textbf{pseudoinverse} $A^{+}$ is turn, divide by each
  singular value, turn \dash{} agreeing with $A^{-1}$ where that exists and
  giving the least-squares answer where it does not. Dividing by the smallest
  $\sigma$ is where the condition number becomes a division you can point at,
  and truncating first is the trade that bounds it.}
\sumitem{33}{\textbf{One factorisation answers six programs' questions}: rank
  is how many $\sigma$ are non-zero, the determinant is their product, the null
  space is spanned by the $v$'s whose $\sigma$ is zero, the best low-rank
  approximation is the first few terms, and $\kappa$ is the first divided by
  the last.}
\end{summarybox}

\canyou

\begin{testexercises}
\item $A$ is $5 \times 3$. Give the number of singular values it has, and say
  whether it can have eigenvalues.
  \answerto{Three \dash{} the smaller of the two dimensions. It cannot have
  eigenvalues at all: $Av$ lives in a space of five dimensions and $v$ in one
  of three, so $Av = \lambda v$ is not a statement that can be made.}
\item $A$ has singular values $6$, $3$ and $1$. Give $\kappa(A)$,
  $\kappa(A\T A)$, and roughly how many decimal digits each costs.
  \answerto{$6$ and $36$. About $\num{0.8}$ of a digit and about $\num{1.6}$ \dash{}
  $\logb{10}6$ and $\logb{10}36$. The point is not the size of either but that
  the second is twice the first, always.}
\item The truncation of $A$ to rank $2$ leaves an error of $1$. What does that
  tell you about $\sigma_3$, and what does it tell you about every other
  rank-$2$ matrix?
  \answerto{$\sigma_3 = 1$, exactly, because the error left is what was
  discarded. And by Eckart--Young no other rank-$2$ matrix is closer, so $1$ is
  not merely this method's error but the best available.}
\item Two matrices have spectra $(100, 99, 98, 97)$ and $(100, 3, 2, 1)$. Say
  which is worth approximating at rank $1$ and why.
  \answerto{The second. Rank $1$ keeps $100^{2}$ out of
  $100^{2}+3^{2}+2^{2}+1^{2}$, which is over $99\%$; the first keeps about a
  quarter, because its four directions matter almost equally. \textbf{Decay is
  the condition, not size.}}
\item You are told a solver returned coefficients that \enquote{look
  reasonable} for a badly conditioned least-squares fit. Say why that is not
  reassuring.
  \answerto{Because a conditioning failure does not produce nonsense. It
  produces plausible numbers of the right sign and magnitude that are wrong in
  the digits nobody checks. \enquote{Looks reasonable} is exactly the symptom.}
\item Why does a matrix with a zero singular value have no inverse, in this
  program's vocabulary?
  \answerto{Because $A^{+}$ divides by each singular value, and dividing by
  zero is where that stops. Equivalently: $\Sigma$ flattens that direction, so
  the map is not one-to-one, which is Program~\ref{prog:P08}'s null space and
  Program~\ref{prog:P09}'s zero determinant.}
\end{testexercises}

\begin{furtherproblems}
\item Show that $\kappa(A) \ge 1$ for every matrix, and say which matrices
  achieve $1$.
  \answerto{$\sigma_{\max} \ge \sigma_{\min}$ by definition, so the ratio is at
  least $1$. Equality means every singular value is the same, so $\Sigma$ is a
  multiple of the identity and $A$ is a rotation scaled by a constant \dash{}
  which stretches every direction alike and so cannot distort a relative
  error at all.}
\item Program~\ref{prog:P10} said a quarter turn has no real eigenvalue. Give
  its singular values.
  \answerto{Both $1$. $A\T A = I$ for any rotation, so every singular value is
  $1$ and $\kappa = 1$. The SVD has no difficulty with the matrix that defeated
  eigenvalues entirely, which is section 1's claim made concrete.}
\item A matrix is multiplied by $1000$. Say what happens to its singular
  values, to its condition number, and to how many digits a solve with it
  loses.
  \answerto{Every singular value is multiplied by $1000$; the condition number
  is unchanged, because it is a ratio; and the digits lost are unchanged. This
  is Program~\ref{prog:P10}'s determinant trap from the other side: scaling
  moves the determinant enormously and the conditioning not at all, which is
  why $\kappa$ is the honest measure of \emph{nearly singular} and $\det$ is
  not.}
\item Why does the pseudoinverse of a matrix with a tiny singular value give
  an answer of enormous length, and what is the standard remedy?
  \answerto{Because it divides by that singular value, so any component of the
  right-hand side along the corresponding direction is amplified by its
  reciprocal. The remedy is to truncate \dash{} drop the offending singular
  values before inverting \dash{} which bounds the amplification at the cost of
  declining to answer in those directions.}
\item Program~\ref{prog:P08} proved that a rank-$r$ bottleneck caps the rank of
  everything after it, and declined to claim that deep stacks lose rank without
  one. Say what measurement would settle the second, and why this book has not
  made it.
  \answerto{The singular values of the representations at each depth: if they
  decay ever more steeply with depth, the rank is falling in the sense the
  claim means. This book has not made it because it needs a trained model,
  which the book does not have and does not download \dash{} and constructing a
  spectrum and reporting it would be a fabrication rather than a measurement.}
\end{furtherproblems}
